As the adoption of modern data-intensive technologies such as artificial intelligence continues to accelerate across industries, the demand for computational power is rising and has been for some time \cite{dhar2020carbon}. Consequently, the electricity use of today's hyperscale data centers is increasing steadily, with 1\% of the global energy consumption being attributed to server farms \cite{masanet2020recalibrating}. With an introduction of carbon pricing mechanisms on the horizon \cite{boyce2018carbon, best2020carbon}, this power demand and the subsequent carbon footprint will become an issue to the operators if their infrastructure is solely powered by the public grid, as a big portion of the energy in today's world is produced using fossil fuels \cite{SustainableEdgeComputing}.

Many cloud providers like Microsoft \footnote{\url{https://blogs.microsoft.com/blog/2020/01/16/microsoft-will-be-carbon-negative-by-2030/}} and Google \footnote{\url{https://www.gstatic.com/gumdrop/sustainability/google-2019-environmental-report.pdf}} have already stated ambitious plans to reduce their environmental impact, which requires them to utilize the promising alternative of distributed renewable energy resources to meet their power demand \cite{MicrogridReview}. In this context, microgrids provide a fitting framework for defining the operation of distributed generation and consumption \cite{Microgrid}, and recent work has outlined the advantages of microgrids over traditional power systems in terms of control, flexibility, economic viability, environmental friendliness, and adaptability \cite{MicrogridPerspectives}. Because of the rather volatile nature of renewable energy sources like solar and wind, the energy output of a localized microgrid will fluctuate over time, which makes the use of energy storage systems such as batteries necessary to deal with upcoming imbalances between demand and supply side \cite{MicrogridStorageAdvances}.

In this complex environment, the task of studying dynamic behaviors like load-, energy-, and battery management becomes increasingly difficult \cite{yu2016distributed}, and the process of researching and developing system configurations, taking these things into account remains challenging due to the limited availability of suitable test environments. Recent work has therefore proposed a tool called Vessim \cite{vessim} to serve as a virtual testbed for carbon-aware systems, where applications have visibility and control over a simulated microgrid. Thereby, co-simulation is utilized, which offers a platform to study coupled heterogeneous systems like microgrids by composing simulations of its different parts \cite{gomes2018co}. This is sensible because at its core, a microgrid comprises an interplay of various components like renewable power sources, energy storage systems, power electronics, and intelligent control algorithms \cite{alzahrani2017modeling}.

Thus far, however, the opportunities of modeling energy storage systems in such testbeds are severely limited. Either only way too simple battery models with nothing more than an adjustable energy level are supported, or one has to work with too complex and slow simulations of a complete underlying microgrid power network, to study accurate physics-based storage models that consider parameters like battery aging and thermal factors. No work has been done, specifying the interfaces that are necessary to integrate energy storage models of different complexity into a microgrid co-simulation, and analyzing, how these various models compare.

\section{Contributions}
The objective of this thesis is to investigate, how different energy storage models can be integrated into a microgrid co-simulation, which can simulate the energy system infrastructure of a data center. These models can range from simple, lossless batteries with the sole purpose of transferring usable energy over the duration of a simulation, to complex models of entire battery packs that take the internal physical processes of batteries as well as the energy losses of the battery packs circuit into account.

This is achieved by extending the already existing co-simulation infrastructure of the Vessim framework \cite{vessim}, which will be described in Section \ref{vessim_infrastructure}, and add more context to specific parts of the design. Thereby, mainly the storage simulation unit is refined and explained in detail so that the described objectives can be met, and the infrastructure can be used as a co-simulation testbed for data centers.

This work then describes the implementation of four different mathematical models of lithium-ion battery packs with battery management systems, using the proposed interface for energy storage units:
\begin{enumerate}
    \item The simplest model (\textit{SimpleBattery}) simulates a basic, lossless battery that has a fixed capacity, and does not consider any inefficiencies or physical properties of a battery pack.
    \item The \textit{CLCBattery} implements the C-L-C model as described by Kazhamiaka et al. \cite{kazhamiaka2019tractable} that is still a simple linear mathematical model for single lithium-ion battery cells, but considers properties such as charging and discharging inefficiencies and power limits.
    \item The \textit{PybammBattery} uses the PyBaMM framework \cite{pybammbatteries} to simulate individual lithium-ion cells that behave according to the Single Particle Model, which takes the electrochemical processes inside a battery into account.
    \item Finally, the \textit{LiionBatteryPack} models the behavior of a full lithium-ion battery pack, where the cells are modeled using PyBaMM's Single Particle Model, and the battery pack's circuit is solved using methods of the Liionpack framework \cite{tranter2022liionpack}.
\end{enumerate}

In the end, the enhanced co-simulation design of Vessim \cite{vessim} is used to analyze, how these different models behave in different scenarios within a simulated microgrid, which can represent a data center's energy system.

\section{Overview}